\section{Triplons and the Kitaev-Heisenberg Model}

In Chapter \ref{chap:kitaev_star}, we considered an extension of the Kitaev honeycomb model to include the more `physical' influence of Heisenberg couplings. This was motivated by research that had shown that a number of candidate materials exist, such as the iridates Na$_2$IrO$_3$ and $\alpha$-Li$_2$IrO$_3$~\cite{singh_antiferromagnetic_2010, singh_relevance_2012,revelli_fingerprints_2020}, as well as $\alpha$-RuCl$_3$~\cite{sandilands_scattering_2015,banerjee_proximate_2016} which have been conjectured to realise a dominant Kitaev interaction, as well as other subdominant terms such as Heisenberg coupling, motivating an effective description in terms of the Kitaev-Heisenberg (KH) model. On the honeycomb, this model was shown to have a ground-state phase diagram containing the expected QSL phases, as well as a variety of classically magnetically ordered phases such as ferro-, antiferromagnetic and stripy ordering~\cite{chaloupka_zigzag_2013,rau_spin-orbit_2016, gohlke_dynamics_2017}. 

In light of the previous chapter, where we showed that the QSL phase was remarkably insensitive to changes in the lattice geometry, we considered how the phase diagram as a whole would depend on lattice geometry. To probe this question, we focussed on the breaking of sublattice symmetry, constructing the KH model on the star lattice -- one of the simplest examples of a lattice with odd cycles. Three methods were used to explore the model. First, we used state-of-the-art DMRG \cite{tenpy} to numerically determine the ground state across the entire parameter range of the KH model. Next, we tested the robustness of the phase diagram using a Majorana mean field calculation, verifying that -- at least in a mean field context -- the phases were robust in the thermodynamic limit. Additionally, we constructed a local transformation, showing that the ground state phase diagram had a hidden symmetry, where different phases could be exactly mapped onto one another by a local spin rotation. We found some familiar phases, the ferromagnetic and QSL phases could be constructed similarly as in honeycomb case. However, the investigation also revealed a pair of valence bond solid (VBS) phases -- a singlet phase and a triplet phase -- which had no honeycomb analogue. These phases formed due to the geometric frustration of N\'eel order. Next, we characterised the singlet VBS phase by constructing a bosonic mean field theory following Sachdev and Bhatt's bond operator formalism \cite{sachdev_bond-operator_1990}. This allowed us to characterise the triplon excitations that form in such a VBS. Finally, by considering the effect of an applied $[1,1,1]$ magnetic field, we showed that the triplon bands realised a rich variety of topological bands and confirmed the existence of topologically protected triplon edge modes. 

The work opens a number of possible avenues of further study. Firstly, we emphasise that the star lattice was chosen because it provided the simplest example of a tri-coordinated lattice breaking bipartite symmetry, and because several studies existed for the star-lattice Heisenberg model, allowing us to benchmark our results~\cite{richter_absence_2004,choy_classification_2009,yang_competing_2010, jahromi_spin-frac12_2018,jahromi_infinite_2018,ran_emergent_2018}. However, the ultimate goal is to understand the properties of the Kitaev-Heisenberg phase diagram on general amorphous lattices. In particular, one motivation for this study was the possibility that, by geometrically frustrating many of the competing phases, it might be possible to tune the lattice geometry such that the QSL phase was stable over a larger region of phase space. This is especially relevant since many candidate Kitaev materials are expected to sit reasonably close to the transition to a QSL phase, but none have been shown conclusively to realise it yet. This question remains open. Certainly in the case of the star lattice the answer is negative (see Fig.~\ref{fig:kh_phase_diagram}), however this does not necessarily apply to all lattices that break sublattice symmetry. Thus, it will be worth repeating this study on a larger range of lattices, and investigating the resulting phase boundaries to see if the robustness of the QSL phase can be improved just through lattice geometry.

There is also a rich variety of further modifications that one could make to the Hamiltonian to further assess the robustness of these phases and their chiral excitations. This is particularly timely in light of recent experimental work which demonstrated the successful synthesis of a spin $1/2$ star lattice structure with antiferromagnetic couplings~\cite{sorolla_synthesis_2020}. The material is expected to have in-triangle couplings much stronger than the inter-triangle couplings. Thus, it would be worth considering how the phase diagram is affected when the relative strength of these two symmetry-inequivalent couplings is adjusted. In addition, it would be interesting to derive microscopic Hamiltonians based on the microscopic orbital configuration and ab-initio calculations, e.g.~considering the effect of including symmetric off-diagonal exchanges~\cite{rau_generic_2014,knolle_dynamics_2018}, as well as Dzyaloshinskii-Moriya interactions \cite{kim_realization_2016, zhang_interplay_2021}.

The new triplet VBS phase we discover around $\varphi \approx 1.7\pi$ in Fig.~\ref{fig:kh_phase_diagram} shows intriguing topological triplet and singlet excitations but many open questions remain. For example, in the absence of an external magnetic field, the dynamics are expected to be identical to the singlet VBS phase due to the four-unit-cell rotation discussed in \textsection\ref{sec:spin_flip_transformation}. However, an applied magnetic field breaks this symmetry, and the possibility remains that this phase could have qualitatively very different properties than the singlet VBS. In addition, it will be important to understand the stability of these topological triplons as well as their experimental signatures.
